\documentclass[11pt,a4paper]{amsart}
\usepackage{amssymb,amsthm,mathtools}
\usepackage[no-math]{fontspec}
\setmonofont{DejaVuSansMono.ttf}[Scale=0.82]
\usepackage[margin=27mm]{geometry}
\usepackage{enumitem}
\usepackage{xcolor}
\usepackage[colorlinks=true,linkcolor=blue!50!black,urlcolor=blue!50!black,citecolor=blue!50!black]{hyperref}

\newtheorem{theorem}{Theorem}
\newtheorem{lemma}[theorem]{Lemma}
\newtheorem{proposition}[theorem]{Proposition}
\theoremstyle{remark}
\newtheorem{remark}[theorem]{Remark}

\newcommand{\R}{\mathbb{R}}
\newcommand{\C}{\mathbb{C}}
\newcommand{\La}{\Lambda}
\newcommand{\lean}[1]{\texttt{#1}}
\renewcommand{\Re}{\operatorname{Re}}
\renewcommand{\Im}{\operatorname{Im}}
\newcommand{\paperdate}{13 November 2026}
\newcommand{\searchdate}{13 November 2026}

\title[Boundedness from below of the $A_4$-symmetric 3HDM, machine-checked]{The conjecture of Lavoura and Milagre on the $A_4$-symmetric three-Higgs-doublet model: machine-checked proofs for $\Lambda_4=0$ and, given the neutral theorem, for $\Lambda_0<0$}

\author{Claude Fable 5}
\thanks{Claude Fable 5 is an AI model made by Anthropic (model ID \texttt{claude-fable-5-2}).
The work was commissioned by Dr.\ Marc Potempa; see the statement on authorship and status in
Section~\ref{sec:status}. Lean sources: \url{https://github.com/MarcPotempa/a4-3hdm-lean-proof}}
\date{\paperdate{} (version 2: Section~\ref{sec:neg} added)}

\emergencystretch=4em
\begin{document}

\begin{abstract}
Lavoura and Milagre (arXiv:2611.05475) derive necessary conditions for the quartic part $V_4$ of the
$A_4$-symmetric potential of three Higgs doublets to be non-negative, and conjecture, on the basis of a
numerical test, that the conditions are also sufficient. We prove the conjecture in the case $\Lambda_4=0$, in
which the symmetry is $S_4$. In this case necessary and sufficient conditions were derived by Ivanov and
Vaz\~ao (JHEP 11 (2020) 104) by a geometric method; we prove in addition that the two sets of conditions are
equivalent for all real couplings, which Lavoura and Milagre had checked numerically. All statements are
formalized in Lean~4 with Mathlib and checked by the Lean kernel. The proof of sufficiency is elementary: it
rests on one inequality for three doublets, which becomes a sum of squares after a gauge rotation.
In addition we prove, for arbitrary $\Lambda_4$ and $\Lambda_0<0$, a duality that reduces non-negativity of
$V_4$ on all field configurations to non-negativity of a second $A_4$ potential on the neutral configurations.
Together with the theorem of Buskin and Ivanov on neutral configurations (which is used as a hypothesis and is
not formalized here) this shows that the conditions of Lavoura and Milagre are sufficient whenever
$\Lambda_0<0$. The case $\Lambda_4\neq0$, $\Lambda_0\ge0$ remains open.
\end{abstract}

\maketitle

\section{Setting and results}

Let $\chi_1,\chi_2,\chi_3\in\C^2$ be three Higgs doublets and $A_{kl}=\chi_k^\dagger\chi_l$. The quartic part
of the $A_4$-symmetric potential is \cite[Eq.~(3)]{LM}
\begin{equation}\label{V4}
\begin{aligned}
V_4={}&\frac{\La_0+\La_3}{3}\bigl(A_{11}^2+A_{22}^2+A_{33}^2\bigr)
+\frac{2\La_0-\La_3}{3}\bigl(A_{11}A_{22}+A_{22}A_{33}+A_{33}A_{11}\bigr)\\
&+\frac{\La_1+\La_2}{2}\bigl(|A_{12}|^2+|A_{23}|^2+|A_{31}|^2\bigr)
+\frac{\La_1-\La_2}{2}\,\Re\bigl(A_{12}^2+A_{23}^2+A_{31}^2\bigr)\\
&+\frac{\La_4}{2}\,\Im\bigl(A_{12}^2+A_{23}^2+A_{31}^2\bigr),
\end{aligned}
\end{equation}
with real couplings $\La_0,\dots,\La_4$. Following \cite{LM}, the potential is called \emph{bounded from below}
if $V_4\ge0$ for all $\chi_1,\chi_2,\chi_3$. For $\La_4=0$ the potential is $S_4$-symmetric.

\subsection*{The conditions of Lavoura and Milagre}
With $R=\sqrt{(\La_1-\La_2)^2+\La_4^2}$, Section~4 of \cite{LM} asks to enforce
\begin{equation}\label{LMa}\tag{43}
\begin{gathered}
\La_0+\La_3\ge0,\qquad \La_0+\La_1\ge0,\qquad \La_0+\frac{\La_1+3\La_2-\sqrt3\,|\La_4|}{4}\ge0,\\
\La_0+\frac{\La_3}{4}+\frac{3(\La_1+\La_2-R)}{8}\ge0,\qquad \La_0+\frac{\La_3}{4}\ge0,
\end{gathered}
\end{equation}
and to exclude each of the following four situations (two inequalities each):
\begin{align}
&|\La_4|-2\sqrt3\,\La_2<0 \ \text{ and }\ \La_0+\frac{4\La_1\La_2-\La_4^2}{16\La_2}<0;\tag{44}\\
&\bigl|3(\La_2-\La_1)\mp\sqrt3\,\La_4\bigr|-3\bigl(3\La_1+\La_2\mp\sqrt3\,\La_4\bigr)<0 \ \text{ and }\
\La_0+\frac{4\La_1\La_2-\La_4^2}{4\bigl(3\La_1+\La_2\mp\sqrt3\,\La_4\bigr)}<0;\tag{45, 46}\\
&\max(0,2\La_0)-\La_3<0 \ \text{ and }\
\max\Bigl(-\La_1-\La_2-2\La_3,\ \La_1+\La_2+\frac{6\La_0\La_3}{\La_0+\La_3}\Bigr)-R<0.\tag{47}
\end{align}
(The upper signs are (45), the lower signs (46); the numbers are those of arXiv:2611.05475v1.)
The authors prove that these conditions are necessary for boundedness from below and write: ``We therefore
conjecture that our necessary conditions for BnessFB are, in fact, necessary and sufficient.''

For $\La_4=0$ they specialise the conditions in their Section~5, Eqs.~(50)--(53): enforce
\begin{equation}\tag{50}
\begin{gathered}
\La_0+\La_3\ge0,\qquad \La_0+\La_1\ge0,\qquad \La_0+\frac{\La_1+3\La_2}{4}\ge0,\\
\La_0+\frac{\La_3}{4}+\frac{3\min(\La_1,\La_2)}{4}\ge0,\qquad \La_0+\frac{\La_3}{4}\ge0,
\end{gathered}
\end{equation}
and exclude
\begin{align}
&\La_2>0 \ \text{ and }\ \La_0+\frac{\La_1}{4}<0;\tag{51}\\
&|\La_2-\La_1|-3\La_1-\La_2<0 \ \text{ and }\ \La_0+\frac{\La_1\La_2}{3\La_1+\La_2}<0;\tag{52}\\
&\max(0,2\La_0)-\La_3<0 \ \text{ and }\
\max\Bigl(-\La_1-\La_2-2\La_3,\ \La_1+\La_2+\frac{6\La_0\La_3}{\La_0+\La_3}\Bigr)-|\La_1-\La_2|<0.\tag{53}
\end{align}

\subsection*{The conditions of Ivanov and Vaz\~ao}
For the $S_4$-symmetric potential, \cite{IV} gives the following necessary and sufficient conditions (Eqs.~(4.5),
(4.8), (4.9) there; quoted as Eqs.~(54)--(56) in \cite{LM}): the eight inequalities
\begin{equation}\tag{54}
\begin{gathered}
\La_0+\La_1\ge0,\quad \La_0+\La_3\ge0,\quad \La_0+\frac{\La_1+3\La_2}{4}\ge0,\quad \La_0+\frac{\La_3+3\La_2}{4}\ge0,\\
\La_0+\frac{\La_1}{4}\ge0,\quad \La_0+\frac{\La_3}{4}\ge0,\quad \La_0+\frac{\La_1+3\La_2}{16}\ge0,\quad
\La_0+\frac{\La_3+3\La_2}{16}\ge0,
\end{gathered}
\end{equation}
and the exclusion of
\begin{align}
&\La_1-|\La_2|>0 \ \text{ and }\ \La_0+\frac{\La_1\La_2}{3\La_1+\La_2}<0,\tag{55}\\
&\La_3-|\La_2|>0 \ \text{ and }\ \La_0+\frac{\La_2\La_3}{\La_2+3\La_3}<0.\tag{56}
\end{align}
About the two sets, \cite{LM} says: ``We have numerically checked that the two methods for guaranteeing BnessFB
are in practice equivalent.''

\subsection*{Results}
The following statements are proved in Lean; the names of the Lean theorems are given in brackets.

\begin{theorem}[\lean{BFB.A4.lm\_conjecture\_S4}]\label{thm:main}
Let $\La_4=0$. Then $V_4\ge0$ for all $\chi_1,\chi_2,\chi_3\in\C^2$ if and only if the conditions
\textup{(43)--(47)} of Lavoura and Milagre hold. That is, their conjecture is true in the $S_4$ case.
\end{theorem}

\begin{theorem}[\lean{BFB.A4.boundedBelow\_iff\_iv}]\label{thm:iv}
Let $\La_4=0$. Then $V_4\ge0$ for all $\chi_1,\chi_2,\chi_3\in\C^2$ if and only if the conditions
\textup{(54)--(56)} of Ivanov and Vaz\~ao hold.
\end{theorem}

\begin{theorem}[\lean{BFB.S4.lmS4\_iff\_iv}]\label{thm:equiv}
For all real $\La_0,\La_1,\La_2,\La_3$, the conditions \textup{(50)--(53)} hold if and only if the conditions
\textup{(54)--(56)} hold.
\end{theorem}

\begin{proposition}[\lean{BFB.A4.gG\_nonneg}]\label{prop:key}
For all $\chi_1,\chi_2,\chi_3\in\C^2$ and every real $\tau\ge\frac34$,
\begin{equation}\label{key}
\tau^2N^2-3\tau E-3(4\tau-3)\,Q\ \ge\ 0,
\end{equation}
where $N=A_{11}+A_{22}+A_{33}$, $E=\sum_{k<l}\bigl(A_{kk}A_{ll}-|A_{kl}|^2\bigr)$ and
$Q=(\Im A_{12})^2+(\Im A_{23})^2+(\Im A_{31})^2$.
\end{proposition}

Theorem~\ref{thm:iv} is the result of Ivanov and Vaz\~ao; what is new here is that it is proved from first
principles (directly from~\eqref{V4}) in a form that a proof assistant can check, by a route different from
the geometric argument of \cite{IV}. Theorem~\ref{thm:equiv} replaces the numerical comparison in \cite{LM} by
a proof. Theorem~\ref{thm:main} follows from Theorems~\ref{thm:iv} and~\ref{thm:equiv} and the fact that
(43)--(47) reduce to (50)--(53) for $\La_4=0$ (\lean{BFB.A4.lm\_iff\_lmS4}).

\begin{remark}[Division by zero]
In Lean, $x/0=0$. This convention does not affect any of the statements: in (44), (45), (46), (52), (55) and
(56) the first inequality of the pair forces the denominator in the second one to be positive; in (47) and (53)
the first inequality gives $\La_3>0$, so the denominator $\La_0+\La_3$ can vanish only if
$\La_0+\La_3/4=-3\La_3/4<0$, in which case the conditions fail already by the last inequality of (43) or (50).
\end{remark}

\section{Proofs}

\subsection{Gram quantities}
Besides $N$, $E$, $Q$ put
\[
Z=\sum_kA_{kk}^2-\sum_{k<l}A_{kk}A_{ll},\qquad P=(\Re A_{12})^2+(\Re A_{23})^2+(\Re A_{31})^2 .
\]
Then $Z\ge0$, $P\ge0$, $Q\ge0$, and $E\ge0$ by the Cauchy--Schwarz inequality. Expanding gives
\begin{equation}\label{rel}
Z+3E+3P+3Q=N^2,\qquad 3V_4=\La_0N^2+\La_3Z+3\La_1P+3\La_2Q\quad(\La_4=0).
\end{equation}

\subsection{The key inequality (Proposition~\ref{prop:key})}
Write $\chi_k=(a_k,b_k)$. The three conditions $\Im(\alpha a_k+\beta b_k)=0$, $k=1,2,3$, are three real linear
equations for the four real unknowns $\Re\alpha,\Im\alpha,\Re\beta,\Im\beta$, so they have a solution
$(\alpha,\beta)\neq(0,0)$ (\lean{exists\_gauge}). The matrix
$\bigl(\begin{smallmatrix}\alpha&\beta\\-\bar\beta&\bar\alpha\end{smallmatrix}\bigr)$ is
$\sqrt{|\alpha|^2+|\beta|^2}$ times a matrix in $SU(2)$; applying it to all three doublets multiplies every
$A_{kl}$ by $s=|\alpha|^2+|\beta|^2>0$ and hence the left-hand side of~\eqref{key} by $s^2$. We may therefore
assume that the upper components are real: $\chi_k=(u_k,d_k)$ with $u_k\in\R$, $d_k\in\C$. With
$m_1=\sum_ku_k^2$ and $m_2=\sum_k|d_k|^2$ one has the polynomial identity (\lean{key\_poly})
\[
2\bigl(\tau^2N^2-3\tau E-3(4\tau-3)Q\bigr)
=\tfrac12\bigl(2\tau m_1-(3-2\tau)m_2\bigr)^2+6\tau\,\Bigl|\sum_ku_kd_k\Bigr|^2
+\tfrac32(4\tau-3)\,\Bigl|\sum_kd_k^2\Bigr|^2 ,
\]
whose right-hand side is non-negative for $\tau\ge\frac34$. \qed

\subsection{Sufficiency in Theorem~\ref{thm:iv}}
Assume (54)--(56) and $\La_4=0$, and let $m=\min(\La_1,\La_3)$. By~\eqref{rel}, since $Z,P\ge0$,
\[
3V_4\ \ge\ \La_0N^2+m\,(Z+3P)+3\La_2Q\ =\ F:=(\La_0+m)N^2-3mE+3(\La_2-m)Q .
\]
It remains to show $F\ge0$. For $m=\La_1$ the conditions (54) and (55) give
\begin{equation}\label{hyp}
\La_0+m\ge0,\qquad \La_0+\frac m4\ge0,\qquad \La_0+\frac{m+3\La_2}{4}\ge0,\qquad
|\La_2|<m\ \Rightarrow\ \La_0(3m+\La_2)+m\La_2\ge0,
\end{equation}
and for $m=\La_3$ the same follows from (54) and (56). We use~\eqref{key}, writing $g(\tau)$ for its left-hand
side, and distinguish three cases (\lean{suff\_aux}).

\emph{Case $m\le\La_2$.} Then $F\ge(\La_0+m)N^2-3mE$. If $m<0$ this is $\ge(\La_0+m)N^2\ge0$. If $m\ge0$,
then $g(\frac34)\ge0$ says $4E\le N^2$, so $F\ge(\La_0+\frac m4)N^2\ge0$.

\emph{Case $m>\La_2$ and $m+\La_2\le0$.} Here $g(\frac32)\ge0$ says $Q\le\frac14N^2-\frac12E$, hence
$F\ge\bigl(\La_0+\frac{m+3\La_2}{4}\bigr)N^2-\frac32(m+\La_2)E\ge0$.

\emph{Case $m>\La_2$ and $m+\La_2>0$.} Then $|\La_2|<m$, $D:=3m+\La_2>0$ and $\tau:=3m/D\ge\frac34$. One
checks the identity
\[
F=\frac D3\,g(\tau)+(\La_0+m-m\tau)\,N^2,\qquad \La_0+m-m\tau=\frac{\La_0(3m+\La_2)+m\La_2}{D}\ \ge0
\]
by the last condition in~\eqref{hyp}. Hence $F\ge0$. \qed

\begin{remark}
The inequalities~\eqref{key} are sharp. On the family $\chi_1=(\omega s,c)$, $\chi_2=(\bar\omega s,c)$,
$\chi_3=(s,c)$ used below ($\omega=e^{2\pi i/3}$, $S=s^2$, $C=c^2$) the left-hand side equals
$\frac94\bigl(2\tau(S+C)-3S\bigr)^2$, and on the family $\chi_1=(0,c)$, $\chi_2=(s,0)$, $\chi_3=(is,0)$ it equals
$\bigl(\tau(C+2S)-3S\bigr)^2$; each vanishes for a suitable ratio $S:C$ whenever $\frac34\le\tau\le\frac32$.
In the plane of $(E/N^2,Q/N^2)$ both families trace the same curve
$\bigl(\sigma(1-\sigma),\sigma^2/4\bigr)$, $0\le\sigma\le1$, and the lines $g(\tau)=0$ are its tangents at
$\sigma=2\tau/3$. This is where the two ``arc'' conditions (55) and (56) come from.
\end{remark}

\begin{remark}
The proof of sufficiency does not use the last two inequalities of (54). Together with the necessity part this
shows that they follow from the remaining conditions (54)--(56).
\end{remark}

\subsection{Necessity in Theorem~\ref{thm:iv}}
Assume $V_4\ge0$ everywhere and $\La_4=0$. Evaluating~\eqref{V4} on two two-parameter families and at two points
gives (\lean{V4\_tfam}, \lean{V4\_qfam}, \lean{V4\_E}, \lean{V4\_E'}), with $S=s^2\ge0$, $C=c^2\ge0$:
\begin{align*}
\chi=\bigl((\omega s,c),(\bar\omega s,c),(s,c)\bigr):&\quad
\tfrac13V_4=\La_0(S+C)^2+\La_1\bigl(C-\tfrac S2\bigr)^2+\tfrac34\La_2S^2,\\
\chi=\bigl((0,c),(s,0),(is,0)\bigr):&\quad
3V_4=\La_0(C+2S)^2+\La_3(C-S)^2+3\La_2S^2,\\
\chi=\bigl((2,0),(-1,\sqrt3),(-1,-\sqrt3)\bigr):&\quad V_4=48\La_0+12\La_1,\\
\chi=\bigl((0,1),(1,0),(0,0)\bigr):&\quad 3V_4=4\La_0+\La_3 .
\end{align*}
The first family at $(S,C)=(0,1),(1,0),(1,1)$ gives the first, third and seventh inequality of (54); the second
family at $(S,C)=(0,1),(1,0),(1,2)$ gives the second, fourth and eighth; the two points give the fifth and the
sixth. If the situation (55) occurred, then $\La_1>|\La_2|$ and the first family at
$(S,C)=(2\La_1,\La_1+\La_2)$ would give
$0\le(3\La_1+\La_2)\bigl(\La_0(3\La_1+\La_2)+\La_1\La_2\bigr)$, contradicting the second inequality of (55).
In the same way the second family at $(S,C)=(\La_3,\La_3+\La_2)$ excludes (56). \qed

\subsection{Equivalence of the two sets of conditions (Theorem~\ref{thm:equiv})}
This is a statement about real numbers only. Both directions are proved by case distinctions on the signs of
the quantities involved and on which argument attains the maxima and minima; each case is closed by linear or
quadratic arithmetic (\lean{lm\_imp\_iv}, \lean{iv\_imp\_lm} in \lean{BFB/S4Equiv.lean}, about 170 lines). For
example, the fifth inequality of (54), $\La_0+\La_1/4\ge0$, follows from the third inequality of (50) together
with the exclusion of (51): if it failed, (51) would force $\La_2\le0$, and then
$\La_0+\La_1/4\ge\La_0+(\La_1+3\La_2)/4\ge0$.

\subsection{From the general conditions to the $S_4$ case}
For $\La_4=0$ one has $R=|\La_1-\La_2|$ and $\La_1+\La_2-R=2\min(\La_1,\La_2)$, so (43) becomes (50); (44)
becomes (51); (45) and (46) both become (52) after division by~3; and (47) becomes (53)
(\lean{BFB.A4.lm\_iff\_lmS4}). With Theorems~\ref{thm:iv} and~\ref{thm:equiv} this proves
Theorem~\ref{thm:main}.


\section{The case $\Lambda_0<0$}\label{sec:neg}

In this section $\La_4$ is arbitrary. Let
\[
M=\begin{pmatrix}\La_1&\La_4/2\\ \La_4/2&\La_2\end{pmatrix},\qquad
m(c)=\La_1(\Re c)^2+\La_2(\Im c)^2+\La_4\,\Re c\,\Im c\quad(c\in\C),
\]
so that, with $T=\sum\Re A_{kl}\Im A_{kl}$ over the pairs $12,23,31$,
\begin{equation}\label{V4m}
3V_4=\La_0N^2+\La_3Z+3\bigl(m(A_{12})+m(A_{23})+m(A_{31})\bigr)=\La_0N^2+\La_3Z+3\La_1P+3\La_2Q+3\La_4T .
\end{equation}
The \emph{dual couplings} are
\[
\La_0'=-\frac1{\La_0},\qquad \La_3'=-\frac4{\La_3},\qquad
\begin{pmatrix}\La_1'&\La_4'/2\\ \La_4'/2&\La_2'\end{pmatrix}=-4M^{-1},
\]
that is $\La_1'=-4\La_2/\delta$, $\La_2'=-4\La_1/\delta$, $\La_4'=4\La_4/\delta$ with
$\delta=\det M=\La_1\La_2-\La_4^2/4$. A configuration is \emph{neutral} if $\chi_k=(v_k,0)$ for $k=1,2,3$.

\begin{theorem}[\lean{BFB.A4.boundedBelow\_of\_dual\_neutral}]\label{thm:dual}
Let $\La_0<0$, $\La_3>0$ and let $M$ be positive definite. If the potential with the dual couplings is
non-negative on all neutral configurations, then $V_4\ge0$ on all field configurations.
\end{theorem}

On neutral configurations the question is settled by the theorem of Buskin and Ivanov \cite[Theorem~1]{BI}. In the
notation of \cite{LM} it says that $V_4\ge0$ on all neutral configurations if and only if the first four
inequalities of (43) hold (footnote in Section~4 of \cite{LM}). We use the ``if'' direction, in the original
formulation of \cite{BI}, as a hypothesis; it is \emph{not} formalized.

\begin{theorem}[\lean{BFB.A4.lm\_sufficient\_of\_L0\_neg}]\label{thm:neg}
Assume the ``if'' direction of Theorem~1 of \cite{BI}. If $\La_0<0$ and the conditions \textup{(43)--(47)} hold, then
$V_4\ge0$ on all field configurations.
\end{theorem}

Since the conditions (43)--(47) are necessary \cite{LM}, Theorem~\ref{thm:neg} says that the conjecture of Lavoura
and Milagre is true for $\La_0<0$, given the theorem of Buskin and Ivanov. (The necessity of the conditions for
$\La_4\neq0$ is proved in \cite{LM}; it is formalized here only for $\La_4=0$.)

\subsection{Proof of Theorem~\ref{thm:dual}}
Three vectors in $\C^2$ are linearly dependent: there is $v\in\C^3$, $v\neq0$, with $v_1\chi_1+v_2\chi_2+v_3\chi_3=0$
(\lean{exists\_null}). Taking the norm of this vector gives (\lean{normSq\_comb})
\[
0=\sum_kp_kA_{kk}+2\,\Re\bigl(w_1A_{12}+w_2A_{23}+w_3A_{31}\bigr),
\]
where $p_k=|v_k|^2$ and $(w_1,w_2,w_3)=(\bar v_1v_2,\ \bar v_2v_3,\ \bar v_3v_1)$.
With $n=p_1+p_2+p_3>0$ the first sum is $\tfrac13nN+\sum_k(p_k-\tfrac n3)(A_{kk}-\tfrac N3)$. The
Cauchy--Schwarz inequality for the Euclidean form on the diagonal part and for the form $m$ on each off-diagonal
entry gives
\[
\Bigl(\sum_k(p_k-\tfrac n3)(A_{kk}-\tfrac N3)\Bigr)^2\le\frac{4\zeta}{9\La_3}\cdot\La_3Z,\qquad
\bigl(2\Re(w_kA_{kl})\bigr)^2\le\frac{4\,m^*(w_k)}{3}\cdot3\,m(A_{kl}),
\]
where $\zeta=\sum_kp_k^2-\sum_{k<l}p_kp_l$ and $m^*(w)=\bigl(\La_2(\Re w)^2+\La_1(\Im w)^2+\La_4\Re w\Im w\bigr)/\delta$
is the form of $M^{-1}$ evaluated at $\bar w$ (\lean{cs\_diag}, \lean{cs\_offdiag}). Adding the four inequalities
(\lean{cs\_add}) yields
\[
\Bigl(\frac{nN}{3}\Bigr)^2\le D\cdot\Bigl(\La_3Z+3\sum m(A_{kl})\Bigr),\qquad
D=\frac{4\zeta}{9\La_3}+\frac43\sum_km^*(w_k).
\]
By~\eqref{V4m}, three times the dual potential on the neutral configuration $(\bar v_1,\bar v_2,\bar v_3)$ equals
$n^2/|\La_0|-9D$ (\lean{V4\_dual\_neutral}); so the hypothesis says $9D\le n^2/|\La_0|$. Hence
$\La_3Z+3\sum m(A_{kl})\ge|\La_0|N^2$, which is $V_4\ge0$. \qed

\subsection{Proof of Theorem~\ref{thm:neg}}
We have to show that (43)--(47) and $\La_0<0$ imply $\La_3>0$, that $M$ is positive definite, and that the dual
couplings satisfy the first four inequalities of (43) (\lean{lm\_dual\_neutralConditions}). Put $\mu=-4\La_0>0$ and,
for $c=x+iy$,
\[
h_0=x,\qquad h_1=-\tfrac x2+\tfrac{\sqrt3}{2}y,\qquad h_2=-\tfrac x2-\tfrac{\sqrt3}{2}y,\qquad h_0+h_1+h_2=0 .
\]
The set $\{\max_kh_k\le\frac12\}$ is the triangle with vertices $-1,-\omega,-\omega^2$. The second and third
inequality of (43) say that $m\ge\mu/4$ at these vertices. On the side $h_0=\frac12$ the form $m$ is a quadratic
polynomial in $y$; the exclusion of (44) says that its minimum, if it lies inside the side, is at least $\mu/4$.
The exclusions of (46) and (45) say the same for the sides $h_1=\frac12$ and $h_2=\frac12$. A quadratic polynomial
on an interval is bounded below by its values at the end points and at an interior minimum
(\lean{quad\_on\_interval}); therefore $m\ge\mu/4$ on the boundary of the triangle, and by homogeneity
$m(c)\ge\mu\,\max_kh_k(c)^2$ for all $c$ (\lean{tri\_max}). Since $m(-c)=m(c)$, it follows that
\[
m(c)\ \ge\ \mu\,h_k(c)^2\qquad\text{for all }c\in\C\text{ and }k=0,1,2
\]
(\lean{tri\_all}). Evaluating this at suitable points gives $\La_1>0$, $\La_2>0$ and
\[
\delta\ge4|\La_0|\La_2,\qquad \delta\ge|\La_0|\,\bigl(3\La_1+\La_2\pm\sqrt3\,\La_4\bigr)
\]
(\lean{tri\_coeff}). In particular $\delta>0$, so $M$ is positive definite, and the two displayed inequalities are
the second and the third inequality of (43) for the dual couplings. The first inequality of (43) gives
$\La_3\ge-\La_0>0$. The first inequality of (43) for the dual couplings is $\La_0+\La_3/4\ge0$, the last
inequality of (43). Finally, since $\La_0<0<\La_3$, the first inequality in (47) holds, so the exclusion of (47)
gives $\La_1+\La_2-R\ge-6\La_0\La_3/(\La_0+\La_3)$, which is equivalent to
$1/|\La_0|\ge1/\La_3+6/(\La_1+\La_2-R)$, the fourth inequality of (43) for the dual couplings
(\lean{dual\_d}). Now Theorem~\ref{thm:dual} and the theorem of Buskin and Ivanov apply. \qed

\subsection{The theorem of Buskin and Ivanov as a hypothesis}
In \cite{BI} the potential is written as
$a\sum_k|z_k|^2+2b\,(|z_1z_2|+|z_2z_3|+|z_3z_1|)+2\Re\bigl(d\,(\bar z_1z_2+\bar z_2z_3+\bar z_3z_1)\bigr)$ for
$z\in\C^3$ (their Eq.~(1)), and Theorem~1 states that it is non-negative on $\C^3$ if and only if $a\ge0$,
$a+b-|d|\ge0$ and $a+2b+2|d|\cos(\gamma+2\pi k/3)\ge0$ for $k=1,2,3$, where $d=|d|e^{i\gamma}$ (their Eqs.~(3), (4)).
The Lean definition \lean{BuskinIvanovTheorem} is the ``if'' direction of this statement, with $\gamma=\arg d$. The
file \lean{BFB/Neutral.lean} proves that for $z_k=v_k^2$,
\[
a=\frac{\La_0+\La_3}{3},\qquad b=\frac{2\La_0-\La_3}{6}+\frac{\La_1+\La_2}{4},\qquad d=\frac{\La_1-\La_2-i\La_4}{4}
\]
this potential is $V_4$ on the neutral configuration $\chi_k=(v_k,0)$ (\lean{V4\_neutral\_eq}), and that the
first four inequalities of (43) imply the conditions of \cite{BI}
(lemma \lean{biConditions\_of\_neutralConditions}).

\section{What is checked, and how}

The Lean files are in the repository named on the first page. \lean{BFB/Defs.lean} contains the definitions that
enter the statements: the potential~\eqref{V4} (\lean{BFB.A4.V4}), boundedness from below, and the three sets of
conditions, transcribed from arXiv:2611.05475v1, the dual couplings, and the statement of the theorem of
Buskin and Ivanov, transcribed from arXiv:2104.11428. The file \lean{Challenge.lean} repeats these definitions and
states the theorems without proof; the program Comparator confirms that the theorems proved in the repository
are exactly those statements and that they depend on no axioms other than the three standard ones of Lean
(\lean{propext}, \lean{Quot.sound}, \lean{Classical.choice}). Whoever wants to rely on the result needs to
read only \lean{Challenge.lean} (about 190 lines) and compare it with the formulas above and in \cite{LM}, \cite{BI}.
Theorem~\ref{thm:neg} is only as good as its hypothesis: should the statement \lean{BuskinIvanovTheorem} deviate
from what is proved in \cite{BI}, the theorem would rest on an unproved assumption. The transcription was
compared with the \LaTeX{} source of arXiv:2104.11428, and the statement was tested numerically.

\section{Relation to earlier work}

The conditions for the $S_4$-symmetric model are due to Ivanov and Vaz\~ao \cite{IV}, who obtained them from a
description of the orbit space. For the $A_4$-symmetric model, conditions for the neutral directions were
conjectured in \cite{IV} and proved by Buskin and Ivanov \cite{BI}; Lavoura and Milagre \cite{LM} added the
charge-breaking directions and stated the conjecture discussed here. A search on \searchdate{} (INSPIRE records
citing \cite{LM} and \cite{IV}, arXiv, GitHub, general web search) found no proof of the conjecture, no proof of
Theorem~\ref{thm:equiv} and no formalization of a three-Higgs-doublet potential in a proof assistant; for the
two-Higgs-doublet model, stability conditions were formalized in Lean by Tooby-Smith \cite{TS}. We also did not
find the duality of Theorem~\ref{thm:dual} in the literature, but the search for it was brief. We may have
missed something.

For $\La_4\neq0$ and $\La_0\ge0$ the conjecture is not proved. Partial results (pen-and-paper arguments, not
formalized and not checked by any human) are collected in the file \lean{notes/general\_case.pdf} of the
repository; Theorem~\ref{thm:dual} is the sufficiency part of ``Theorem~1'' of that note.

\section{Authorship and status}\label{sec:status}

The proof, its Lean formalization and this description were produced by Claude Fable 5, an AI model made by
Anthropic (model ID \texttt{claude-fable-5-2}, working in Claude Code), on 11--13 November 2026. The work was
commissioned by Dr.\ Marc Potempa, who set the task and decided to publish it, but did not contribute to the
mathematical content. The authors of \cite{LM} were informed before publication.

The Lean proof was checked with the standard Lean kernel (Lean~4.27.0, Mathlib v4.27.0), through
\lean{lake build} and Comparator. No human has read the proof or this description. Where this description and
the Lean files differ, the Lean files are authoritative. The transcription of the conditions from \cite{LM} into
Lean is the part that the kernel cannot check; it was compared with the \LaTeX{} source of arXiv:2611.05475v1 and,
for \cite{IV}, with the journal version, but not by a second person.

\begin{thebibliography}{9}
\bibitem{LM} L.~Lavoura, A.~Milagre, \emph{Complete set of bounded-from-below conditions for the
$A_4$-symmetric three-Higgs-doublet model}, arXiv:2611.05475v1 (3 November 2026).
\bibitem{IV} I.~P.~Ivanov, F.~Vaz\~ao, \emph{Yet another lesson on the stability conditions in multi-Higgs
potentials}, JHEP 11 (2020) 104, arXiv:2006.00036.
\bibitem{BI} N.~Buskin, I.~P.~Ivanov, \emph{Bounded-from-below conditions for $A_4$-symmetric 3HDM},
J.~Phys.~A 54 (2021) 325401, arXiv:2104.11428.
\bibitem{TS} J.~Tooby-Smith, \emph{Formalizing the stability of the two Higgs doublet model potential into
Lean: identifying an error in the literature}, arXiv:2603.08139.
\end{thebibliography}

\end{document}
